%=====================================================================
\chapter{Enough PHP, SQL and Twig to Read the Shop}\label{app:languages}
%=====================================================================
\normalfigures

The sibling courses give this appendix to mathematics. This one spends it on
language, because that is what this course actually assumes and does not teach.

You are not asked to be a strong PHP programmer. Until \chref{modules} you will
read far more code than you write, and this appendix is calibrated to that:
enough to follow what the platform is doing, and enough to change it when
\chref{modules} asks you to.

\section{PHP, in Two Pages}

\subsection{The shape of a file}

PHP is a language for producing a response to an HTTP request. A file runs top
to bottom, every time the page is requested, and whatever it prints becomes the
page.

\begin{lstlisting}[language=php]
<?php
declare(strict_types=1);

$name  = 'Keyboard';       // string, single quotes: no interpolation
$price = 2499.00;          // float
$qty   = 3;                // int
$total = $price * $qty;

// Double quotes interpolate; single quotes do not. This is the single
// most common source of confusion in a first week of PHP.
echo "Total for $name: $total";   // Total for Keyboard: 7497
echo 'Total for $name';           // Total for $name
\end{lstlisting}

\begin{alertbox}[Three things that surprise everyone]
\textbf{Variables start with \texttt{\$}, always} --- including inside strings
and in function parameters.

\smallskip
\textbf{\texttt{==} is not \texttt{===}.} The two-character form converts types
before comparing, so \texttt{'0' == false} is true and \texttt{'abc' == 0} was
true in older versions. \textbf{Use \texttt{===} unless you have a reason not
to}, and in \chref{security} you will see what the other one costs.

\smallskip
\textbf{\texttt{declare(strict\_types=1)} is not the default.} Without it PHP
quietly converts \texttt{"5 apples"} into \texttt{5} when a function wants an
integer. OpenCart's own code does not use it; your modules should.
\end{alertbox}

\subsection{Arrays, which are really two things}

\begin{lstlisting}[language=php]
<?php
// A list.
$sizes = ['S', 'M', 'L'];
echo $sizes[0];                      // S
echo count($sizes);                  // 3

// A map -- the same type, with keys. Almost every piece of data in
// OpenCart arrives as one of these.
$product = [
    'product_id' => 42,
    'name'       => 'Keyboard',
    'price'      => 2499.00,
];
echo $product['name'];               // Keyboard

// Walking one. Note the "as $key => $value" form.
foreach ($product as $field => $value) {
    echo "$field = $value\n";
}

// The ?? operator: use this value, or that one if it is absent.
// Without it, a missing key is a warning and a null.
$image = $product['image'] ?? 'placeholder.png';
\end{lstlisting}

\subsection{Functions, classes and the arrow}

\begin{lstlisting}[language=php]
<?php
function formatPrice(float $amount, string $symbol = 'Rs. '): string
{
    return $symbol . number_format($amount, 2);
}

class ShippingQuote
{
    private float $rate;

    public function __construct(float $rate)
    {
        $this->rate = $rate;         // $this-> reaches a property
    }

    public function forWeight(float $kg): float
    {
        return $this->rate * max(1.0, $kg);
    }
}

$q = new ShippingQuote(120.0);
echo formatPrice($q->forWeight(2.5));    // Rs. 300.00
\end{lstlisting}

\begin{conceptbox}[The three arrows, which are different]
PHP uses three symbols that beginners merge into one.

\begin{tight}
  \item \texttt{\$obj->method()} --- call a method on an object, or read a
        property. An \emph{object}.
  \item \texttt{\$array['key']} --- read a key from an array. An \emph{array}.
        Square brackets, quoted key.
  \item \texttt{Class::CONSTANT} --- reach something on the class itself rather
        than on an instance.
\end{tight}

\smallskip
OpenCart hands you arrays far more often than objects, so \texttt{['key']} is
the one you will type most. \texttt{\$this->request->get['product\_id']} is all
three ideas in one expression: a property of a property, then an array key.
\end{conceptbox}

\subsection{Where request data arrives, and why it is dangerous}

\begin{lstlisting}[language=php]
<?php
// Submitted by a form with method="get"  -> $_GET
// Submitted by a form with method="post" -> $_POST
$qty = $_POST['quantity'] ?? 1;

// EVERY value in those arrays is a STRING, and every one of them was
// typed by somebody who may not be a customer. Chapter 16 is about
// what happens when this is forgotten; Chapter 15 is about sessions,
// which is how the shop remembers who you are between requests.
$qty = (int)$qty;                    // now it is certainly an integer
if ($qty < 1 || $qty > 999) {
    $qty = 1;                        // and certainly in range
}
\end{lstlisting}

\section{SQL, as This Course Uses It}

You met SQL in \emph{Database Systems}. This section is a reminder of the five
statements this course needs and the two ideas it leans on.

\begin{lstlisting}[language=sql]
-- Read. The only one you will type often.
SELECT p.product_id, pd.name, p.price
  FROM oc_product p
  JOIN oc_product_description pd ON pd.product_id = p.product_id
 WHERE pd.language_id = 1
   AND p.status = 1
 ORDER BY p.price DESC
 LIMIT 20;

-- Count. Chapter 6 shows why a category page runs one of these as
-- well as the query above, every single time.
SELECT COUNT(*) FROM oc_product WHERE status = 1;

-- Change one row. The WHERE is not optional; without it you have
-- just repriced the entire catalogue.
UPDATE oc_product SET price = 2999.00 WHERE product_id = 42;

-- Remove. Same warning, more so.
DELETE FROM oc_product_to_category WHERE product_id = 42;

-- Make a query fast, or fail to -- see Chapter 6.
CREATE INDEX idx_sort_order ON oc_product (sort_order);
\end{lstlisting}

\begin{definitionbox}[The two ideas everything else rests on]
\term{A join} combines rows from two tables on a matching column. OpenCart
splits a product across eighteen tables (\chref{catalogue}), so almost every
useful question needs at least one join.

\smallskip
\term{An index} is a sorted structure the database can search instead of
reading every row. It makes reads faster and writes slower, and it is useless
if the query hides the column inside a function --- \texttt{ORDER BY
LCASE(name)} cannot use an index on \texttt{name}, which is the central finding
of \chref{catalogue}.
\end{definitionbox}

\begin{examplebox}[EXPLAIN, which answers "why is this slow?"]
Put \texttt{EXPLAIN} in front of any \texttt{SELECT} and the database tells you
what it intends to do rather than doing it. Three words in the output carry
nearly all the information:

\begin{tight}
  \item \textbf{\texttt{type: ALL}} --- a full table scan. Every row read. At
        fifty thousand products this is the one to fix.
  \item \textbf{\texttt{Using filesort}} --- the result had to be sorted after
        it was fetched, because no index gave it in the right order.
  \item \textbf{\texttt{Using temporary}} --- a temporary table was built to
        hold intermediate results, usually because of \texttt{DISTINCT} or
        \texttt{GROUP BY}.
\end{tight}

\smallskip
\chref{catalogue} runs \texttt{EXPLAIN} against the real category query and
finds all three at once.
\end{examplebox}

\section{Twig, Which Is Just the Page}

OpenCart's templates are Twig. A template is HTML with three kinds of hole in
it, and that is very nearly the whole language.

\begin{lstlisting}[language=twig]
{# A comment. It does not reach the browser. #}

{# 1. Print something. #}
<h1>{{ heading_title }}</h1>
<span class="price">{{ product.price }}</span>

{# 2. Do something -- a loop or a condition. #}
{% if products %}
  <ul>
    {% for product in products %}
      <li>
        <a href="{{ product.href }}">{{ product.name }}</a>
        {% if product.special %}
          <s>{{ product.price }}</s> {{ product.special }}
        {% else %}
          {{ product.price }}
        {% endif %}
      </li>
    {% endfor %}
  </ul>
{% else %}
  <p>{{ text_empty }}</p>
{% endif %}

{# 3. Filter -- transform a value on the way out. #}
<p>{{ product.description|striptags|slice(0, 120) }}</p>
<p>{{ product.name|escape }}</p>
\end{lstlisting}

\begin{alertbox}[\texttt{|escape} and what it prevents]
Twig escapes output by default in most configurations, and that default is the
single thing standing between a product description written by somebody else
and a cross-site scripting hole in your shop.

\smallskip
Twig also offers \texttt{|raw}, which turns the escaping off. There are
legitimate uses, and every one of them is a decision to trust the source of
that value completely. \chref{security} looks for places where somebody has
used it without meaning to.
\end{alertbox}

\section{YAML, in Ten Lines}

The Compose file is YAML, and YAML has exactly two rules worth knowing.

\begin{lstlisting}[language=yaml]
# 1. Indentation is structure. Two spaces per level, NEVER tabs --
#    a tab is a syntax error and the message will not say so clearly.
services:
  db:
    image: mariadb:11
    environment:
      MARIADB_DATABASE: opencart

# 2. A leading dash is a list item.
    ports:
      - "8090:80"
      - "3306:3306"
\end{lstlisting}

That is enough to read \texttt{code/docker-compose.yml} and to change the one
thing \chref{stack} asks you to change, which is a port number.

\section{Where To Look Next}

\rowcolors{2}{white}{lightbg}
\begin{longtable}{@{}L{3.2cm}L{11.8cm}@{}}
\toprule
\rowcolor{primary!15}
\textbf{For} & \textbf{Go to} \\
\midrule
\endhead
PHP & \texttt{php.net/manual} --- the official manual, and unusually good. Each
function page has user-contributed notes; treat those as hints, not authority \\
SQL & The MariaDB knowledge base, and your \emph{Database Systems} notes for
the theory this appendix skips \\
Twig & \texttt{twig.symfony.com/doc} --- the ``for template designers'' page is
all you need \\
OpenCart & \texttt{docs.opencart.com}, and the source in your own container,
which is more current than the documentation \\
\bottomrule
\end{longtable}

\begin{conceptbox}[Read the source; it is right there]
The shop's entire source is inside your container, and you can read it:

\smallskip
\texttt{docker compose exec shop \textbackslash}\\
\texttt{\phantom{xx}cat /var/www/html/catalog/model/catalog/product.php}

\smallskip
This is not a last resort. \chref{catalogue} finds the category page's real
query this way, and it is the only way to answer ``what does this page actually
do?'' with certainty rather than with a plausible guess. The documentation
describes a version; the container holds yours.
\end{conceptbox}
